\section{Backward demands and parameter-changing induction}
\label{sec:demand}

\subsection{Why ordinary uniform precision is wasteful}

Many constructions transform an input approximation into an output
approximation.  To establish the output at precision \(k\), a local step may
need its predecessor at a different precision \(D(n,k)\).  Repeated products,
ODE steps, root refinements, matrix iterations, and error-budget pipelines all
have this shape.  A naive proof chooses one large precision valid for every
intermediate stage.  This often introduces unnecessary monotonicity hypotheses,
and it can place an all-precision assertion inside the induction formula.

The correct finite object is a backward demand schedule.

\subsection{The schedule}

Fix a terminal time \(N\), requested output precision \(k\), and a
primitive-recursive local demand function \(D(n,r)\).  Define
\[
 h_{N,k}(N)=k,
 \qquad
 h_{N,k}(t)=D(t,h_{N,k}(t+1))\quad(t<N).
\]
For primitive recursion it is convenient to compute backward distance:
\begin{align*}
 b(0)&=k,\\
 b(j+1)&=D(N-j-1,b(j)),\\
 h_{N,k}(t)&=b(N-t).
\end{align*}
Truncated subtraction handles the total numerical definition; the theorem is
used only under \(t\le N\).

\begin{lemma}[Demand equations]
\label{lem:demand-equations}
iPRA proves
\[
 h_{N,k}(N)=k
 \quad\text{and}\quad
 t<N\to h_{N,k}(t)=D(t,h_{N,k}(t+1)).
\]
It also proves the horizon coherence identity
\[
 h_{N+1,k}(t)=h_{N,D(N,k)}(t)\qquad(t\le N).
\]
\end{lemma}

\begin{proof}
The final equation is the zero recursion equation for \(b\).  For the step,
\(N-t=(N-(t+1))+1\) under \(t<N\), so one unfolding of the recursion gives the
claim.  Horizon coherence is induction on the backward distance; the numerical
indices agree by elementary truncated-subtraction identities.  All statements
are equations between fixed primitive-recursive terms.
\end{proof}

\subsection{Demand-passing transfer}

Let \(P(n,r)\) be a bounded property saying that the stage-\(n\) object is
correct at requested observation \(r\).  Assume
\begin{align*}
 P(0,h_{N,k}(0)),\tag{Base}\label{eq:demand-base}\\
 P(n,D(n,r))\to P(n+1,r)\qquad(n<N).\tag{Step}\label{eq:demand-step}
\end{align*}

\begin{theorem}[Finite demand-passing transfer]
\label{thm:demand-transfer}
From \eqref{eq:demand-base} and \eqref{eq:demand-step}, iPRA proves
\(P(N,k)\).  No monotonicity of \(D\) or of \(P\) is required.
\end{theorem}

\begin{proof}
Induct on \(t\le N\) with the Booleanized quantifier-free invariant
\[
  P(t,h_{N,k}(t)).
\]
The base is \eqref{eq:demand-base}.  At a successor \(t+1\), the schedule
equation rewrites the induction hypothesis as
\(P(t,D(t,h_{N,k}(t+1)))\).  Apply \eqref{eq:demand-step} with
\(r=h_{N,k}(t+1)\).  At \(t=N\), the final schedule equation rewrites the
conclusion to \(P(N,k)\).
\end{proof}

The theorem is more than a semantic induction trick.  When \(P\) has been
Booleanized and \(D\) is represented by a fixed \(\PR\) symbol, the proof is
an instance of the checked \code{primitiveRecursionInvariant} transformer.
The schedule is a first primitive recursion; the forward computation and proof
are a second primitive recursion.  This compiles parameter-changing recursion
into ordinary iPRA recursion without adding an induction schema.

\subsection{Variable finite predecessor sets}

Some outputs need several predecessor observations.  Let
\(q(n,j,r)\), for \(r<c(n,j)\), list all predecessor queries needed to prove
observation \(j\) at stage \(n+1\).  Define the finite envelope
\[
 E(n,h)=\max_{j\le h}\max_{r<c(n,j)} q(n,j,r).
\]
If the presentation provides all observations up to a prefix bound, the local
step can use \(E(n,h)\) as one predecessor demand.

\begin{theorem}[Finite-query induction]
\label{thm:finite-query-induction}
Suppose every base observation \(A(0,j)\) holds and
\[
 \bigl(\forall r<c(n,j)\;A(n,q(n,j,r))\bigr)\to A(n+1,j).
\]
Then iPRA proves \(A(N,j)\) for every \(N,j\), provided the bounded universal
premise is compiled to a Boolean fold.
\end{theorem}

\begin{proof}
Let \(P(n,h)\) mean that \(A(n,j)\) holds for every \(j\le h\).  The envelope
ensures that every query needed for an output \(j\le h\) is covered by
\(P(n,E(n,h))\).  Apply \cref{thm:demand-transfer} with demand \(E\), and
recover the desired component by bounded elimination after the induction.
\end{proof}

\subsection{Adaptive finite queries}

A query may depend on earlier finite comparison results.  There are two sound
compilation strategies.

\paragraph{Stateful strategy.}
Include the query log in the trace state.  The invariant says that every entry
recorded so far is an instance of the global input contract.  At a successor,
the next query is a primitive-recursive term in the current state.  Universal
hypotheses are instantiated at this term, and a Boolean fold updates the log
validity bit.

\paragraph{Certificate strategy.}
An external producer supplies a query log and answer table.  A fixed checker
verifies that every query follows the program, every answer satisfies the
input contract at the recorded precision, and every state transition is valid.
The object proof then reasons only about this finite table.

Neither route places ``for every possible future precision'' inside the
induction invariant.  The requested final precision is fixed during the finite
proof and universally introduced afterward.

\subsection{Demand and error ledgers}

A useful implementation record attaches to every certified operation:
\begin{enumerate}[label=(\roman*)]
\item a backward demand transformer \(D_F(n)\);
\item a finite error bound \(e_F(n)\);
\item a list of guards that must hold at the demanded input precision; and
\item an object proof that these data imply the output contract.
\end{enumerate}
For a composite \(G\circ F\), the demand is
\(D_F(D_G(n))\), while the error ledger is obtained by substituting the
intermediate magnitude and guard bounds.  Thus precision bookkeeping becomes
an explicit compiler pass instead of ad hoc arithmetic in every theorem.
